\scanpage{328}\mspage{244}
% l. 1–5
Dans le cas où $\mathscr D$ est de la forme $C\times J$, alors $\unsure{\mathscr N|\mathscr M}$ est de la \unsure{formé|forme} \deleted{$\unsure{\operatorname{Ind}}$} $M\times J$, avec $M=\operatorname{Ind}_{\unsure{\pi|\kappa}}(C)$, et l'inclusion (4.10.3) n'est autre \deleted{que} (à isom près) que
\[
 \underline{\mathrm{Hom}}(J,C)\longrightarrow\underline{\mathrm{Hom}}(J,M).
\]

% l. 6–11
\unsure{Revenant} au cas général, on voit encore que $1^\circ)$ Le deuxième \unsure{mb|mbr} de (4.10.3) est stable par $\underline{\lim}$ grandes devant $\pi$, lesquelles se calculent fibre \unsure{à|:} fibre, et $2^\circ)$ Le premier membre de (4.10.3) est formé d'objets $\pi$-accessibles du second.

% l. 12–14
D'où un foncteur \uline{pleinement} fidèle
\[
 \operatorname{Ind}_{\pi}\bigl(\underline{\Gamma}(\mathscr D/J)\bigr)\hookrightarrow
 \underline{\Gamma}(\mathscr N/J). \tag*{(4.10.4)}
\]

% l. 15–19
\uline{Théorème} 4.10.1. Supposons qu'on ait
\[
 \operatorname{Card}\unsure{\mathrm{red}|\mathrm{réd}}\,J\leq\pi, \tag*{(4.10.1.1)}
\]
et qu'il existe un ordinal limite $\alpha$, tel que \deleted{1) pour tout ordinal $\beta<\alpha$, on ait} \uline{1)} $\operatorname{Card}\deleted{\beta}\alpha\leq\pi$, et \uline{2)} les limites inductives \unsure{de|des}
% archive sheet number: 327
